\chapter{Tablero Kanban: listado completo de tareas}\label{cap:AnexoJ}

Este anexo recoge, en la Tabla~\ref{tab:kanban}, el listado completo de tareas que componen el tablero Kanban del proyecto, organizadas por épica. Los identificadores de tarea (\texttt{T-XX}) son los empleados en el tablero durante el desarrollo para referenciar cada ítem. La última columna recoge la prioridad de atención dentro de cada épica o, en su caso, el estado de exclusión: \textbf{Alta} (bloquea el avance de otras tareas), \textbf{Media} (necesaria pero no bloqueante), \textbf{Baja} (mejora o calidad sin impacto crítico en el camino principal) y \textbf{Descartada} (identificada durante el desarrollo pero excluida explícitamente del alcance del proyecto).

\begin{center}
\small
\begin{longtable}{llp{7cm}c}
\caption{Listado completo de tareas del tablero Kanban con su prioridad o estado, agrupadas por épica.}\label{tab:kanban}\\
\hline
\textbf{ID} & \textbf{Épica} & \textbf{Tarea} & \textbf{Prioridad / estado} \\
\hline
\endfirsthead
\hline
\textbf{ID} & \textbf{Épica} & \textbf{Tarea} & \textbf{Prioridad / estado} \\
\hline
\endhead
\hline
\endfoot

T-01 & Datos & Descarga del corpus CodiEsp (diagnósticos, procedimientos y variante extendida) & Alta \\
T-02 & Datos & Análisis estadístico del corpus: distribución de longitudes, frecuencia de códigos y desequilibrio de clases & Alta \\
T-03 & Datos & Desarrollo del script de \emph{scraping} sobre el portal eCIE-maps para obtener el catálogo CIE-10-ES & Alta \\
T-04 & Datos & Análisis estadístico del catálogo CIE-10-ES: cobertura, distribución por capítulos y vocabulario & Media \\
T-05 & Datos & Generación de los ficheros CSV de entrenamiento, validación y test & Alta \\
T-48 & Datos & Aumentación del corpus CodiEsp mediante retrotraducción con Azure Translator (F0, nivel gratuito) usando tres pivotes (EN, DE, FR), cuadruplicando el corpus de 500 a 2.000 documentos (evaluada y finalmente descartada por no superar la referencia de 500 documentos; véase el Anexo~\ref{cap:AnexoL}) & Media \\
\hline
T-06 & Modelo & Comparativa de modelos base candidatos (IIC/RigoBERTa-Clinical, mmBERT, BETO) & Alta \\
T-07 & Modelo & Implementación del clasificador plano multilabel basado en BERT & Alta \\
T-08 & Modelo & Configuración del pipeline de entrenamiento (BCE con ponderación, AdamW, \emph{scheduler}) & Alta \\
T-09 & Modelo & Implementación del clasificador de diccionario como \emph{baseline} & Media \\
T-10 & Modelo & Entrenamiento y selección del mejor \emph{checkpoint} por F1-micro en validación & Alta \\
T-11 & Modelo & Evaluación sobre el conjunto de test con métricas de precisión, \emph{recall} y F1 & Media \\
T-46 & Modelo & Activación de FlashAttention~2 como implementación eficiente del mecanismo de atención: integrada en la imagen base GPU (\texttt{cie10-ml-base}) con fallback a SDPA; activa en todas las ejecuciones de entrenamiento & Media \\
\hline
T-12 & AI Engine & Implementación del servicio FastAPI con entorno virtualizado & Alta \\
T-13 & AI Engine & Endpoint \texttt{POST /predict} y \texttt{GET /} & Alta \\
T-14 & AI Engine & Desarrollo del servicio \emph{mock} para entornos CI y desarrollo sin GPU & Alta \\
T-15 & AI Engine & Dockerización con imagen base optimizada para PyTorch & Media \\
\hline
T-16 & Backend & Configuración del entorno y herramientas de calidad (linting, análisis estático) & Alta \\
T-17 & Backend & Sistema de autenticación: registro, confirmación por email e inicio de sesión & Alta \\
T-18 & Backend & Arquitectura CQRS/Event Sourcing con Commanded y PostgreSQL como \emph{event store} & Alta \\
T-19 & Backend & Canal WebSocket para comunicación en tiempo real con el frontend & Alta \\
T-20 & Backend & Pipeline de análisis: recepción del informe, llamada al AI Engine y almacenamiento & Alta \\
T-21 & Backend & Soporte de internacionalización (i18n) en cinco idiomas & Media \\
T-22 & Backend & Suite de pruebas con ExUnit e informe de cobertura con excoveralls & Media \\
T-23 & Backend & Dockerización del entorno de desarrollo, pruebas y producción & Media \\
\hline
T-24 & Frontend & Página de autenticación con formularios de registro e inicio de sesión & Alta \\
T-25 & Frontend & Interfaz de chat con renderizado de tarjetas tipadas (resumen y códigos) & Alta \\
T-26 & Frontend & Panel lateral con historial de análisis y controles de sesión & Alta \\
T-27 & Frontend & Componentes de validación y rechazo de códigos predichos con motivo & Alta \\
T-28 & Frontend & Selector de idioma y carga de traducciones desde el backend & Media \\
T-29 & Frontend & Configuración centralizada de URLs mediante variables de entorno & Baja \\
T-30 & Frontend & Dockerización del entorno de desarrollo, pruebas y producción & Media \\
\hline
T-31 & Infraestructura & Configuración de Docker Compose para todos los servicios & Alta \\
T-32 & Infraestructura & Makefile con \emph{targets} para el ciclo de vida completo del proyecto & Media \\
T-33 & Infraestructura & Configuración de commitlint y Husky para validación de mensajes de commit & Baja \\
T-34 & Infraestructura & Workflow de GitHub Actions para integración continua & Media \\
T-35 & Infraestructura & Dockerización del entorno de entrenamiento con soporte CPU y GPU & Alta \\
T-36 & Infraestructura & Despliegue con Docker Compose y Traefik como proxy inverso (TLS automático de Let's Encrypt), publicado mediante un túnel de Cloudflare & Media \\
T-37 & Infraestructura & Monitorización con Prometheus y Grafana: el backend expone métricas de latencia HTTP, queries Ecto y VM BEAM en \texttt{:9568/metrics} mediante \texttt{telemetry\_metrics\_prometheus}; Prometheus scrapea todos los servicios y Grafana los visualiza con datasource provisionado automáticamente & Baja \\
T-38 & Infraestructura & Gestión de errores y excepciones con Sentry: captura de excepciones HTTP mediante \texttt{Sentry.PlugCapture} en el endpoint Phoenix, captura de crashes OTP (GenServer, comandos Commanded) mediante \texttt{Sentry.LoggerHandler}, con stacktrace completo y etiqueta de entorno configurada vía \texttt{SENTRY\_DSN} & Baja \\
T-47 & Infraestructura & Configuración del servidor de email en producción con registros SPF, DKIM y política DMARC para garantizar la entregabilidad y autenticidad de los emails transaccionales & Media \\
\hline
T-39 & Modelo & Bucle de aprendizaje activo: reentrenamiento incremental con códigos validados por el codificador & Descartada \\
T-40 & Modelo & Soporte de codificación de procedimientos CIE-10-ES (ICD-10-PCS) además de diagnósticos & Descartada \\
T-41 & Backend & Control de acceso por roles: administrador, codificador clínico y supervisor & Descartada \\
T-42 & Backend & Integración con el estándar HL7 FHIR para importación de informes desde sistemas HIS & Descartada \\
T-43 & Backend & Autenticación federada mediante OAuth~2.0 / SSO hospitalario & Descartada \\
T-44 & Frontend & Procesamiento por lotes: envío y codificación de múltiples informes de forma simultánea & Descartada \\
T-45 & Modelo & Motor de codificación por búsqueda inversa: indexación del catálogo CIE-10-ES con BM25/TF-IDF y recuperación de códigos por similitud léxica con el informe & Descartada \\
\hline
\end{longtable}
\end{center}
